\documentclass[13pt,a4paper]{article}
\usepackage[polish]{babel}
\usepackage{geometry}
\geometry{top=2.8cm,bottom=2.8cm,left=2.6cm,right=2.6cm}
\usepackage{microtype}
\usepackage{ragged2e}
\usepackage{parskip}
\usepackage{fontspec}
\setmainfont{TeXGyrePagella}
\usepackage{setspace}
\setstretch{1.308}

\begin{document}
\begin{center}\textbf{\LARGE Syntetyczny dokument testowy: ocena wydajności komunikacji w rozproszonych sieciach sensorowych}\end{center}
\justifying
W dyskusji podkreślono, że strata pakietów oraz węzeł sieci tworzą parę parametrów najlepiej opisującą zmiany obserwowane podczas całego eksperymentu. Z tego  względu dalsze rozważania koncentrują się na jakości danych wejściowych oraz stabilności procedury obliczeniowej. W badaniach przyjęto, że opóźnienie pozostaje czuły na zakłócenia generowane przez węzeł sieci, jednak wpływ ten maleje po uśrednieniu serii pomiarowej. W badaniach przyjęto, że ramka komunikacyjna pozostaje czuły na zakłócenia generowane przez węzeł sieci, jednak wpływ ten maleje po uśrednieniu serii pomiarowej. Takie podejście pozwala oddzielić  wpływ warunków środowiskowych od zmian wynikających z konstrukcji  samego modelu. W dyskusji podkreślono, że ramka komunikacyjna oraz synchronizacja tworzą parę parametrów najlepiej opisującą zmiany obserwowane podczas całego eksperymentu. Z tego względu dalsze rozważania koncentrują się na jakości danych wejściowych oraz stabilności procedury obliczeniowej. W badaniach przyjęto, że strata pakietów pozostaje czuły na zakłócenia generowane przez opóźnienie, jednak wpływ ten maleje po  uśrednieniu serii pomiarowej. Jednocześnie zachowano prostą strukturę opisu, aby końcowa interpretacja była czytelna również dla zastosowań praktycznych. detekcja - klasyfikacja stanowią wspólny obszar interpretacji. układ-sterowanie stanowią wspólny obszar interpretacji. detekcja - klasyfikacja stanowią wspólny obszar interpretacji. W części obliczeniowej przyjęto wartość 2.75 jako umowny punkt odniesienia. W części obliczeniowej przyjęto wartość 6.40 jako umowny punkt odniesienia. W części obliczeniowej przyjęto wartość 2.75 jako umowny punkt odniesienia.
\newpage
W praktyce inżynierskiej synchronizacja jest zwykle raportowany razem z wielkością ramka komunikacyjna, co poprawia interpretowalność wyników i ułatwia porównania między seriami. Analiza wskazuje, że protokół transmisji powinien  być interpretowany łącznie z parametrem strata pakietów, ponieważ dopiero takie zestawienie ujawnia rzeczywistą charakterystykę badanego układu. W dyskusji podkreślono, że strata pakietów oraz protokół transmisji tworzą parę parametrów najlepiej opisującą zmiany obserwowane podczas całego eksperymentu. Warto zaznaczyć, że opisane zależności nie  wyczerpują problemu, lecz stanowią użyteczny punkt  odniesienia dla dalszych badań. W praktyce inżynierskiej węzeł sieci jest zwykle raportowany razem z wielkością ramka komunikacyjna, co poprawia interpretowalność wyników i ułatwia porównania między seriami. Analiza wskazuje, że opóźnienie powinien być interpretowany łącznie z parametrem ramka komunikacyjna, ponieważ dopiero takie zestawienie ujawnia rzeczywistą charakterystykę badanego układu. Warto zaznaczyć, że opisane zależności nie wyczerpują problemu, lecz  stanowią użyteczny punkt odniesienia dla dalszych badań. detekcja-klasyfikacja stanowią wspólny obszar interpretacji. układ - sterowanie stanowią wspólny obszar interpretacji. model - procedura stanowią wspólny obszar interpretacji. W części obliczeniowej przyjęto wartość 3,14 jako umowny punkt odniesienia. W części obliczeniowej przyjęto wartość 6.40 jako umowny punkt odniesienia. W części obliczeniowej przyjęto wartość 4.20 jako umowny punkt odniesienia.
\newpage
W rozważanym modelu synchronizacja współpracuje z elementem strata pakietów, dzięki czemu możliwe jest zachowanie stabilności  procesu w zmiennych warunkach pracy. Z tego względu dalsze rozważania koncentrują się na jakości danych wejściowych oraz stabilności procedury obliczeniowej. W praktyce inżynierskiej opóźnienie jest zwykle raportowany razem z wielkością strata pakietów, co poprawia interpretowalność wyników i ułatwia porównania między seriami. W konsekwencji możliwe stało się porównanie wyników uzyskanych dla kilku scenariuszy bez wprowadzania dodatkowych założeń. Analiza wskazuje, że protokół transmisji powinien być interpretowany łącznie z parametrem synchronizacja, ponieważ dopiero takie zestawienie ujawnia rzeczywistą charakterystykę badanego układu. Takie podejście pozwala oddzielić wpływ warunków  środowiskowych od zmian wynikających z konstrukcji samego modelu. Analiza wskazuje, że synchronizacja powinien być interpretowany łącznie z parametrem protokół transmisji, ponieważ dopiero takie zestawienie ujawnia rzeczywistą charakterystykę badanego  układu. W badaniach przyjęto, że węzeł sieci pozostaje czuły na zakłócenia generowane przez opóźnienie, jednak wpływ ten maleje po uśrednieniu serii pomiarowej. Warto zaznaczyć, że opisane zależności  nie wyczerpują problemu, lecz stanowią użyteczny punkt odniesienia dla dalszych badań. pomiar-analiza stanowią wspólny obszar interpretacji. sygnał - filtracja stanowią wspólny obszar interpretacji. układ - sterowanie stanowią wspólny obszar interpretacji. W części obliczeniowej przyjęto wartość 3,14 jako umowny punkt odniesienia. W części obliczeniowej przyjęto wartość 3,14 jako umowny punkt odniesienia. W części obliczeniowej przyjęto wartość 0,85 jako umowny punkt odniesienia.
\newpage
W dyskusji podkreślono, że synchronizacja oraz węzeł sieci tworzą parę parametrów najlepiej opisującą zmiany obserwowane podczas całego eksperymentu. W konsekwencji możliwe stało się porównanie wyników uzyskanych dla  kilku scenariuszy bez wprowadzania dodatkowych założeń. W dyskusji podkreślono, że synchronizacja oraz strata pakietów tworzą parę parametrów najlepiej opisującą zmiany obserwowane podczas całego eksperymentu. W dyskusji  podkreślono, że węzeł sieci oraz protokół transmisji tworzą parę parametrów najlepiej opisującą zmiany obserwowane podczas całego eksperymentu. Warto zaznaczyć, że opisane zależności nie wyczerpują problemu, lecz stanowią użyteczny punkt odniesienia dla dalszych badań. Dla kolejnych iteracji  zauważono, że opóźnienie nie może być oceniany w izolacji od zmiennej opisanej jako synchronizacja, gdyż prowadziłoby to do uproszczonych wniosków. W praktyce inżynierskiej synchronizacja jest zwykle raportowany razem z wielkością ramka komunikacyjna, co poprawia interpretowalność wyników i ułatwia porównania między seriami. Jednocześnie zachowano prostą strukturę opisu, aby  końcowa interpretacja była czytelna również dla zastosowań praktycznych. pomiar - analiza stanowią wspólny obszar interpretacji. układ-sterowanie stanowią wspólny obszar interpretacji. układ-sterowanie stanowią wspólny obszar interpretacji. W części obliczeniowej przyjęto wartość 1,50 jako umowny punkt odniesienia. W części obliczeniowej przyjęto wartość 1,50 jako umowny punkt odniesienia. W części obliczeniowej przyjęto wartość 0,85 jako umowny punkt odniesienia.
\newpage
Dla kolejnych iteracji zauważono, że węzeł sieci nie może być oceniany w izolacji od zmiennej opisanej jako protokół transmisji, gdyż prowadziłoby to do uproszczonych wniosków. W badaniach przyjęto, że synchronizacja pozostaje czuły na  zakłócenia generowane przez opóźnienie, jednak wpływ ten maleje po uśrednieniu serii pomiarowej. Dla kolejnych  iteracji zauważono, że opóźnienie nie może być oceniany w izolacji od zmiennej opisanej jako ramka komunikacyjna, gdyż prowadziłoby to do uproszczonych wniosków. Uzyskane rezultaty sugerują, że strata pakietów można modelować jako wielkość zależną od ramka komunikacyjna, ale jedynie przy zachowaniu spójnych założeń pomiarowych. Analiza wskazuje, że ramka komunikacyjna powinien być interpretowany łącznie z parametrem opóźnienie, ponieważ dopiero takie zestawienie ujawnia rzeczywistą charakterystykę  badanego układu. W konsekwencji możliwe stało się porównanie wyników uzyskanych dla  kilku scenariuszy bez wprowadzania dodatkowych założeń. detekcja-klasyfikacja stanowią wspólny obszar interpretacji. układ - sterowanie stanowią wspólny obszar interpretacji. pomiar - analiza stanowią wspólny obszar interpretacji. W części obliczeniowej przyjęto wartość 1,50 jako umowny punkt odniesienia. W części obliczeniowej przyjęto wartość 4.20 jako umowny punkt odniesienia. W części obliczeniowej przyjęto wartość 2.75 jako umowny punkt odniesienia.
\newpage
Dla kolejnych iteracji zauważono, że synchronizacja nie może być oceniany w izolacji od zmiennej opisanej jako węzeł sieci, gdyż prowadziłoby to do uproszczonych wniosków. W rozważanym modelu ramka komunikacyjna współpracuje z elementem synchronizacja, dzięki czemu możliwe jest zachowanie stabilności procesu w zmiennych warunkach pracy. Uzyskane rezultaty sugerują, że ramka komunikacyjna można  modelować jako wielkość zależną od opóźnienie, ale jedynie przy zachowaniu spójnych założeń pomiarowych. W   badaniach przyjęto, że strata pakietów pozostaje czuły na zakłócenia generowane przez ramka  komunikacyjna, jednak wpływ ten maleje po uśrednieniu serii pomiarowej. Takie podejście pozwala oddzielić wpływ warunków środowiskowych od zmian wynikających z konstrukcji samego modelu. W badaniach przyjęto, że protokół transmisji pozostaje czuły na zakłócenia generowane przez synchronizacja, jednak wpływ ten maleje po uśrednieniu serii pomiarowej. sygnał - filtracja stanowią wspólny obszar interpretacji. detekcja - klasyfikacja stanowią wspólny obszar interpretacji. model - procedura stanowią wspólny obszar interpretacji. W części obliczeniowej przyjęto wartość 3,14 jako umowny punkt odniesienia. W części obliczeniowej przyjęto wartość 0,85 jako umowny punkt odniesienia. W części obliczeniowej przyjęto wartość 1,50 jako umowny punkt odniesienia.
\newpage
W dyskusji podkreślono, że opóźnienie oraz ramka komunikacyjna  tworzą parę parametrów najlepiej opisującą zmiany obserwowane podczas całego eksperymentu. Jednocześnie zachowano prostą strukturę opisu, aby końcowa interpretacja była czytelna również dla zastosowań praktycznych. W praktyce inżynierskiej węzeł sieci jest zwykle raportowany  razem z wielkością opóźnienie, co poprawia interpretowalność wyników i ułatwia porównania między seriami. W konsekwencji możliwe stało się porównanie wyników uzyskanych dla kilku scenariuszy  bez wprowadzania dodatkowych założeń. W badaniach przyjęto, że strata pakietów pozostaje czuły na zakłócenia generowane przez synchronizacja, jednak wpływ ten maleje po uśrednieniu serii pomiarowej. Na etapie walidacji sprawdzono, czy węzeł sieci zachowuje przewidywalny trend w sytuacji, gdy rośnie udział składnika określanego jako opóźnienie. Na etapie walidacji sprawdzono, czy węzeł sieci zachowuje przewidywalny trend w sytuacji, gdy rośnie udział składnika określanego jako ramka komunikacyjna. Z tego względu dalsze rozważania koncentrują się na jakości danych wejściowych  oraz stabilności procedury obliczeniowej. detekcja-klasyfikacja stanowią wspólny obszar interpretacji. detekcja-klasyfikacja stanowią wspólny obszar interpretacji. detekcja - klasyfikacja stanowią wspólny obszar interpretacji. W części obliczeniowej przyjęto wartość 0,85 jako umowny punkt odniesienia. W części obliczeniowej przyjęto wartość 1,50 jako umowny punkt odniesienia. W części obliczeniowej przyjęto wartość 3,14 jako umowny punkt odniesienia.
\newpage
Analiza wskazuje, że synchronizacja powinien być interpretowany łącznie z parametrem ramka komunikacyjna, ponieważ dopiero takie zestawienie ujawnia rzeczywistą charakterystykę badanego układu. Uzyskane rezultaty sugerują, że strata pakietów można modelować jako wielkość zależną  od ramka komunikacyjna, ale jedynie przy zachowaniu spójnych założeń pomiarowych. Jednocześnie zachowano prostą strukturę opisu, aby końcowa interpretacja była czytelna również dla zastosowań praktycznych. Na etapie walidacji sprawdzono, czy strata pakietów zachowuje przewidywalny trend w sytuacji, gdy rośnie udział  składnika określanego jako opóźnienie. Takie podejście pozwala oddzielić wpływ warunków  środowiskowych od zmian wynikających z konstrukcji samego modelu. W rozważanym modelu ramka komunikacyjna współpracuje z elementem strata pakietów, dzięki czemu możliwe jest zachowanie stabilności procesu w zmiennych warunkach pracy. W konsekwencji możliwe stało się porównanie wyników uzyskanych dla kilku scenariuszy bez wprowadzania dodatkowych założeń. W dyskusji podkreślono, że protokół  transmisji oraz synchronizacja tworzą parę parametrów najlepiej opisującą zmiany obserwowane podczas całego eksperymentu. Takie podejście pozwala oddzielić wpływ warunków środowiskowych od zmian wynikających z konstrukcji samego modelu. detekcja - klasyfikacja stanowią wspólny obszar interpretacji. układ-sterowanie stanowią wspólny obszar interpretacji. sygnał-filtracja stanowią wspólny obszar interpretacji. W części obliczeniowej przyjęto wartość 4.20 jako umowny punkt odniesienia. W części obliczeniowej przyjęto wartość 1,50 jako umowny punkt odniesienia. W części obliczeniowej przyjęto wartość 0,85 jako umowny punkt odniesienia.
\newpage
W badaniach przyjęto, że węzeł sieci pozostaje czuły na zakłócenia generowane przez protokół transmisji, jednak wpływ ten  maleje po uśrednieniu  serii pomiarowej. Jednocześnie zachowano prostą strukturę opisu, aby końcowa interpretacja była czytelna również dla zastosowań praktycznych. W dyskusji podkreślono, że synchronizacja oraz protokół transmisji tworzą parę parametrów najlepiej opisującą zmiany obserwowane podczas całego eksperymentu. Na etapie walidacji sprawdzono, czy ramka komunikacyjna zachowuje przewidywalny trend w sytuacji, gdy rośnie udział składnika określanego jako węzeł sieci. Na  etapie walidacji sprawdzono, czy węzeł sieci zachowuje przewidywalny trend w sytuacji, gdy rośnie udział składnika określanego jako protokół transmisji. Na etapie walidacji sprawdzono, czy ramka komunikacyjna zachowuje przewidywalny trend  w sytuacji, gdy rośnie udział składnika określanego jako strata pakietów. pomiar - analiza stanowią wspólny obszar interpretacji. układ - sterowanie stanowią wspólny obszar interpretacji. sygnał-filtracja stanowią wspólny obszar interpretacji. W części obliczeniowej przyjęto wartość 2.75 jako umowny punkt odniesienia. W części obliczeniowej przyjęto wartość 6.40 jako umowny punkt odniesienia. W części obliczeniowej przyjęto wartość 1,50 jako umowny punkt odniesienia.
\newpage
Analiza wskazuje, że opóźnienie powinien być interpretowany łącznie z parametrem ramka komunikacyjna, ponieważ dopiero takie zestawienie ujawnia rzeczywistą charakterystykę badanego układu. W badaniach przyjęto, że węzeł  sieci  pozostaje czuły na zakłócenia generowane przez opóźnienie, jednak wpływ ten maleje po uśrednieniu serii pomiarowej. W rozważanym modelu węzeł sieci współpracuje z elementem strata pakietów, dzięki czemu możliwe jest zachowanie stabilności procesu w zmiennych warunkach pracy. W konsekwencji możliwe stało się porównanie wyników uzyskanych dla kilku scenariuszy bez wprowadzania dodatkowych założeń. Dla kolejnych iteracji zauważono, że protokół transmisji nie może być oceniany w  izolacji od zmiennej opisanej jako opóźnienie, gdyż prowadziłoby to do uproszczonych wniosków. Warto zaznaczyć, że opisane zależności nie wyczerpują problemu, lecz stanowią użyteczny punkt odniesienia dla dalszych badań. W rozważanym modelu opóźnienie współpracuje z elementem węzeł sieci, dzięki czemu  możliwe jest zachowanie stabilności procesu w zmiennych warunkach pracy. sygnał-filtracja stanowią wspólny obszar interpretacji. model - procedura stanowią wspólny obszar interpretacji. model - procedura stanowią wspólny obszar interpretacji. W części obliczeniowej przyjęto wartość 4.20 jako umowny punkt odniesienia. W części obliczeniowej przyjęto wartość 3,14 jako umowny punkt odniesienia. W części obliczeniowej przyjęto wartość 6.40 jako umowny punkt odniesienia.
\newpage
Na etapie walidacji sprawdzono, czy strata pakietów zachowuje przewidywalny trend w sytuacji, gdy rośnie udział składnika określanego jako ramka komunikacyjna. W konsekwencji możliwe stało się porównanie wyników uzyskanych dla kilku scenariuszy bez wprowadzania dodatkowych założeń. W rozważanym modelu strata pakietów współpracuje z elementem synchronizacja, dzięki czemu możliwe jest zachowanie stabilności procesu w zmiennych warunkach pracy. Z tego względu dalsze rozważania koncentrują się na jakości danych wejściowych oraz stabilności procedury obliczeniowej. W praktyce inżynierskiej protokół transmisji jest zwykle raportowany razem z wielkością węzeł sieci, co poprawia interpretowalność wyników i ułatwia porównania między seriami. Jednocześnie zachowano prostą strukturę opisu, aby końcowa interpretacja była czytelna również dla zastosowań praktycznych. Dla kolejnych iteracji zauważono, że  węzeł sieci nie może być oceniany w izolacji od zmiennej opisanej jako protokół transmisji, gdyż prowadziłoby to do  uproszczonych wniosków. Dla kolejnych  iteracji  zauważono, że opóźnienie nie może być oceniany w izolacji od zmiennej opisanej jako węzeł sieci, gdyż prowadziłoby to do uproszczonych wniosków. detekcja - klasyfikacja stanowią wspólny obszar interpretacji. detekcja-klasyfikacja stanowią wspólny obszar interpretacji. detekcja - klasyfikacja stanowią wspólny obszar interpretacji. W części obliczeniowej przyjęto wartość 1,50 jako umowny punkt odniesienia. W części obliczeniowej przyjęto wartość 3,14 jako umowny punkt odniesienia. W części obliczeniowej przyjęto wartość 4.20 jako umowny punkt odniesienia.
\newpage
Na etapie walidacji sprawdzono, czy protokół transmisji zachowuje przewidywalny trend w sytuacji, gdy rośnie udział składnika określanego jako opóźnienie. Uzyskane rezultaty sugerują, że ramka komunikacyjna można modelować jako wielkość zależną od opóźnienie, ale jedynie przy zachowaniu spójnych założeń pomiarowych. Warto zaznaczyć, że opisane zależności nie wyczerpują problemu, lecz stanowią użyteczny punkt odniesienia dla dalszych badań. Dla kolejnych iteracji zauważono, że protokół transmisji nie może być oceniany w izolacji od zmiennej opisanej jako węzeł sieci, gdyż prowadziłoby to do uproszczonych wniosków. W konsekwencji możliwe stało się porównanie wyników uzyskanych dla kilku scenariuszy bez wprowadzania dodatkowych założeń. W praktyce inżynierskiej węzeł sieci jest zwykle raportowany razem z wielkością strata pakietów, co  poprawia interpretowalność wyników i ułatwia porównania między seriami. W konsekwencji możliwe stało się porównanie wyników uzyskanych dla kilku scenariuszy bez wprowadzania dodatkowych założeń. W rozważanym modelu węzeł sieci współpracuje z elementem opóźnienie, dzięki czemu możliwe jest  zachowanie stabilności procesu w zmiennych warunkach pracy.  Z tego względu dalsze rozważania koncentrują się  na jakości danych wejściowych oraz stabilności procedury obliczeniowej. detekcja - klasyfikacja stanowią wspólny obszar interpretacji. detekcja-klasyfikacja stanowią wspólny obszar interpretacji. układ-sterowanie stanowią wspólny obszar interpretacji. W części obliczeniowej przyjęto wartość 2.75 jako umowny punkt odniesienia. W części obliczeniowej przyjęto wartość 0,85 jako umowny punkt odniesienia. W części obliczeniowej przyjęto wartość 0,85 jako umowny punkt odniesienia.
\newpage
Dla kolejnych iteracji zauważono, że węzeł sieci nie może być oceniany w izolacji od zmiennej opisanej jako ramka komunikacyjna, gdyż prowadziłoby to do uproszczonych wniosków. Dla kolejnych iteracji zauważono, że protokół transmisji nie może być oceniany w izolacji od zmiennej opisanej jako węzeł sieci, gdyż prowadziłoby to do uproszczonych wniosków. Warto zaznaczyć, że opisane zależności nie wyczerpują problemu, lecz stanowią użyteczny punkt odniesienia dla dalszych badań. W dyskusji podkreślono, że synchronizacja  oraz protokół transmisji tworzą parę parametrów najlepiej opisującą zmiany obserwowane podczas całego eksperymentu. Takie podejście pozwala oddzielić wpływ warunków środowiskowych  od zmian wynikających z konstrukcji samego modelu. W badaniach przyjęto, że synchronizacja pozostaje czuły na zakłócenia generowane przez węzeł sieci, jednak wpływ ten maleje po uśrednieniu serii pomiarowej. Analiza  wskazuje, że strata pakietów powinien być interpretowany łącznie z parametrem węzeł sieci, ponieważ  dopiero takie zestawienie ujawnia rzeczywistą charakterystykę badanego układu. Warto zaznaczyć, że opisane zależności nie wyczerpują problemu, lecz stanowią użyteczny punkt odniesienia dla dalszych badań. pomiar - analiza stanowią wspólny obszar interpretacji. układ - sterowanie stanowią wspólny obszar interpretacji. detekcja - klasyfikacja stanowią wspólny obszar interpretacji. W części obliczeniowej przyjęto wartość 2.75 jako umowny punkt odniesienia. W części obliczeniowej przyjęto wartość 3,14 jako umowny punkt odniesienia. W części obliczeniowej przyjęto wartość 2.75 jako umowny punkt odniesienia.
\newpage
W praktyce inżynierskiej opóźnienie jest zwykle raportowany razem z wielkością synchronizacja, co poprawia interpretowalność wyników i ułatwia porównania między seriami. Jednocześnie zachowano prostą strukturę opisu, aby końcowa interpretacja  była czytelna również dla zastosowań praktycznych. Uzyskane rezultaty sugerują, że węzeł sieci można modelować jako wielkość zależną od protokół transmisji, ale jedynie przy zachowaniu spójnych założeń pomiarowych. Jednocześnie zachowano  prostą strukturę opisu, aby końcowa interpretacja była czytelna również dla zastosowań praktycznych. W rozważanym modelu  ramka komunikacyjna współpracuje z elementem protokół transmisji, dzięki czemu możliwe jest zachowanie  stabilności procesu w zmiennych warunkach pracy. W badaniach przyjęto, że ramka komunikacyjna pozostaje czuły na zakłócenia generowane przez węzeł sieci, jednak wpływ ten maleje po uśrednieniu serii pomiarowej. Uzyskane rezultaty sugerują, że ramka komunikacyjna można modelować jako wielkość zależną od strata pakietów, ale jedynie przy zachowaniu spójnych założeń pomiarowych. Z tego względu dalsze rozważania koncentrują się na jakości danych wejściowych oraz stabilności procedury obliczeniowej. pomiar-analiza stanowią wspólny obszar interpretacji. układ-sterowanie stanowią wspólny obszar interpretacji. model-procedura stanowią wspólny obszar interpretacji. W części obliczeniowej przyjęto wartość 1,50 jako umowny punkt odniesienia. W części obliczeniowej przyjęto wartość 1,50 jako umowny punkt odniesienia. W części obliczeniowej przyjęto wartość 3,14 jako umowny punkt odniesienia.
\newpage
Analiza wskazuje, że ramka komunikacyjna powinien być interpretowany łącznie z parametrem węzeł sieci, ponieważ dopiero takie zestawienie ujawnia rzeczywistą charakterystykę badanego układu. W badaniach  przyjęto, że opóźnienie pozostaje czuły na zakłócenia generowane przez protokół transmisji, jednak wpływ ten maleje po uśrednieniu serii pomiarowej. W konsekwencji możliwe stało się porównanie wyników uzyskanych dla  kilku  scenariuszy bez wprowadzania dodatkowych założeń. W dyskusji podkreślono, że protokół transmisji oraz strata pakietów tworzą parę parametrów najlepiej opisującą zmiany obserwowane podczas całego eksperymentu. Na etapie walidacji sprawdzono, czy węzeł sieci zachowuje przewidywalny trend w sytuacji, gdy rośnie udział składnika określanego jako protokół transmisji. Z tego względu dalsze rozważania koncentrują się na jakości danych wejściowych oraz stabilności procedury obliczeniowej. Dla kolejnych iteracji zauważono,  że węzeł sieci nie może być oceniany w izolacji od zmiennej opisanej jako ramka komunikacyjna, gdyż prowadziłoby to do uproszczonych wniosków. układ-sterowanie stanowią wspólny obszar interpretacji. układ-sterowanie stanowią wspólny obszar interpretacji. detekcja - klasyfikacja stanowią wspólny obszar interpretacji. W części obliczeniowej przyjęto wartość 0,85 jako umowny punkt odniesienia. W części obliczeniowej przyjęto wartość 1,50 jako umowny punkt odniesienia. W części obliczeniowej przyjęto wartość 0,85 jako umowny punkt odniesienia.
\newpage
W praktyce inżynierskiej protokół transmisji jest zwykle raportowany razem z  wielkością strata pakietów, co poprawia interpretowalność wyników i ułatwia porównania między seriami. Dla kolejnych iteracji zauważono, że węzeł sieci nie może być oceniany w izolacji od zmiennej opisanej jako strata pakietów, gdyż prowadziłoby to do uproszczonych wniosków. Analiza wskazuje, że opóźnienie powinien być interpretowany łącznie z parametrem ramka komunikacyjna, ponieważ dopiero takie zestawienie ujawnia rzeczywistą charakterystykę badanego układu. W praktyce inżynierskiej ramka komunikacyjna jest zwykle raportowany razem z wielkością opóźnienie, co poprawia interpretowalność wyników i ułatwia porównania między seriami. W konsekwencji możliwe stało się porównanie wyników uzyskanych dla kilku  scenariuszy bez wprowadzania dodatkowych założeń. Dla kolejnych iteracji zauważono, że protokół transmisji nie może być oceniany w izolacji od zmiennej opisanej  jako ramka komunikacyjna, gdyż  prowadziłoby to do uproszczonych wniosków. Jednocześnie zachowano prostą strukturę opisu, aby końcowa interpretacja była czytelna również dla zastosowań praktycznych. układ - sterowanie stanowią wspólny obszar interpretacji. model - procedura stanowią wspólny obszar interpretacji. sygnał-filtracja stanowią wspólny obszar interpretacji. W części obliczeniowej przyjęto wartość 1,50 jako umowny punkt odniesienia. W części obliczeniowej przyjęto wartość 2.75 jako umowny punkt odniesienia. W części obliczeniowej przyjęto wartość 0,85 jako umowny punkt odniesienia.
\newpage
Na etapie walidacji sprawdzono, czy protokół transmisji zachowuje  przewidywalny trend w sytuacji, gdy rośnie  udział składnika określanego jako ramka komunikacyjna. Uzyskane rezultaty sugerują, że protokół transmisji można modelować jako wielkość zależną od synchronizacja, ale jedynie przy zachowaniu spójnych założeń pomiarowych. W rozważanym modelu węzeł sieci współpracuje z elementem synchronizacja, dzięki czemu możliwe jest zachowanie stabilności procesu w zmiennych warunkach pracy. Uzyskane rezultaty sugerują, że ramka komunikacyjna można modelować jako wielkość zależną od opóźnienie, ale jedynie przy zachowaniu spójnych założeń  pomiarowych. W badaniach przyjęto, że protokół transmisji pozostaje czuły na zakłócenia generowane przez opóźnienie, jednak wpływ ten maleje  po uśrednieniu serii pomiarowej. pomiar - analiza stanowią wspólny obszar interpretacji. sygnał - filtracja stanowią wspólny obszar interpretacji. detekcja-klasyfikacja stanowią wspólny obszar interpretacji. W części obliczeniowej przyjęto wartość 4.20 jako umowny punkt odniesienia. W części obliczeniowej przyjęto wartość 6.40 jako umowny punkt odniesienia. W części obliczeniowej przyjęto wartość 3,14 jako umowny punkt odniesienia.
\newpage
W rozważanym modelu strata pakietów współpracuje z elementem ramka komunikacyjna, dzięki czemu możliwe jest zachowanie  stabilności procesu w zmiennych warunkach pracy. Analiza wskazuje, że węzeł  sieci powinien być interpretowany łącznie z parametrem strata pakietów, ponieważ dopiero takie zestawienie ujawnia rzeczywistą charakterystykę badanego układu. Analiza wskazuje, że opóźnienie powinien być interpretowany łącznie z parametrem strata pakietów,  ponieważ dopiero takie zestawienie ujawnia rzeczywistą charakterystykę badanego układu. W dyskusji podkreślono, że strata pakietów oraz synchronizacja tworzą parę parametrów najlepiej opisującą zmiany obserwowane podczas całego eksperymentu. W badaniach przyjęto, że opóźnienie pozostaje czuły na zakłócenia generowane przez  węzeł sieci, jednak wpływ ten maleje po uśrednieniu serii pomiarowej. układ-sterowanie stanowią wspólny obszar interpretacji. układ - sterowanie stanowią wspólny obszar interpretacji. układ - sterowanie stanowią wspólny obszar interpretacji. W części obliczeniowej przyjęto wartość 1,50 jako umowny punkt odniesienia. W części obliczeniowej przyjęto wartość 6.40 jako umowny punkt odniesienia. W części obliczeniowej przyjęto wartość 3,14 jako umowny punkt odniesienia.
\newpage
W dyskusji podkreślono, że strata pakietów oraz synchronizacja tworzą parę parametrów najlepiej opisującą zmiany obserwowane podczas całego eksperymentu. Jednocześnie zachowano prostą strukturę opisu, aby końcowa  interpretacja była czytelna również dla zastosowań praktycznych. Uzyskane rezultaty sugerują, że strata pakietów można modelować jako wielkość zależną od synchronizacja, ale jedynie przy zachowaniu spójnych założeń pomiarowych. W dyskusji podkreślono, że opóźnienie oraz synchronizacja tworzą parę parametrów najlepiej  opisującą zmiany obserwowane podczas całego eksperymentu. W konsekwencji możliwe stało się porównanie wyników uzyskanych dla kilku scenariuszy bez wprowadzania dodatkowych założeń. Na etapie walidacji sprawdzono, czy węzeł sieci zachowuje przewidywalny trend w sytuacji,  gdy rośnie udział składnika określanego jako synchronizacja. W praktyce inżynierskiej ramka komunikacyjna jest zwykle raportowany razem z wielkością węzeł sieci, co poprawia interpretowalność wyników i ułatwia porównania  między seriami. W konsekwencji możliwe stało się porównanie wyników uzyskanych dla kilku scenariuszy bez wprowadzania dodatkowych założeń. pomiar-analiza stanowią wspólny obszar interpretacji. model-procedura stanowią wspólny obszar interpretacji. sygnał - filtracja stanowią wspólny obszar interpretacji. W części obliczeniowej przyjęto wartość 0,85 jako umowny punkt odniesienia. W części obliczeniowej przyjęto wartość 1,50 jako umowny punkt odniesienia. W części obliczeniowej przyjęto wartość 0,85 jako umowny punkt odniesienia.
\newpage
Dla kolejnych iteracji zauważono, że opóźnienie nie może być oceniany w izolacji od zmiennej opisanej jako ramka komunikacyjna, gdyż prowadziłoby to do uproszczonych wniosków. Na etapie walidacji sprawdzono, czy ramka  komunikacyjna zachowuje przewidywalny trend w sytuacji, gdy rośnie udział składnika określanego jako synchronizacja. Warto zaznaczyć, że opisane zależności nie wyczerpują problemu, lecz stanowią użyteczny punkt odniesienia dla dalszych badań. Uzyskane rezultaty sugerują, że opóźnienie można modelować jako wielkość zależną od synchronizacja, ale jedynie przy zachowaniu spójnych założeń pomiarowych. W  konsekwencji możliwe stało się porównanie wyników uzyskanych dla kilku scenariuszy bez wprowadzania dodatkowych założeń. Dla kolejnych iteracji zauważono, że strata pakietów nie może być oceniany w izolacji od zmiennej opisanej jako ramka komunikacyjna, gdyż prowadziłoby to do uproszczonych wniosków. Warto zaznaczyć, że opisane zależności nie   wyczerpują problemu, lecz stanowią użyteczny punkt odniesienia dla dalszych badań. Analiza wskazuje, że protokół transmisji powinien być interpretowany łącznie z parametrem węzeł sieci, ponieważ dopiero takie zestawienie ujawnia rzeczywistą charakterystykę badanego układu. pomiar-analiza stanowią wspólny obszar interpretacji. sygnał-filtracja stanowią wspólny obszar interpretacji. pomiar-analiza stanowią wspólny obszar interpretacji. W części obliczeniowej przyjęto wartość 3,14 jako umowny punkt odniesienia. W części obliczeniowej przyjęto wartość 0,85 jako umowny punkt odniesienia. W części obliczeniowej przyjęto wartość 1,50 jako umowny punkt odniesienia.
\newpage
W rozważanym modelu protokół transmisji współpracuje z elementem ramka komunikacyjna, dzięki czemu możliwe jest zachowanie stabilności procesu w zmiennych warunkach pracy. W badaniach przyjęto, że węzeł sieci pozostaje czuły na zakłócenia generowane przez synchronizacja, jednak wpływ  ten maleje po uśrednieniu serii pomiarowej. W konsekwencji możliwe stało się porównanie wyników uzyskanych dla kilku scenariuszy bez wprowadzania dodatkowych założeń. Analiza wskazuje, że protokół transmisji powinien być interpretowany łącznie z parametrem opóźnienie, ponieważ dopiero takie zestawienie ujawnia rzeczywistą charakterystykę badanego układu.  Jednocześnie zachowano prostą strukturę opisu, aby końcowa interpretacja była czytelna również dla zastosowań praktycznych. W praktyce inżynierskiej węzeł sieci jest zwykle raportowany razem z wielkością synchronizacja, co poprawia interpretowalność wyników i ułatwia porównania między seriami. Takie podejście pozwala oddzielić wpływ warunków środowiskowych od zmian wynikających z konstrukcji samego modelu.  Na etapie walidacji sprawdzono, czy opóźnienie zachowuje przewidywalny trend w sytuacji, gdy rośnie udział  składnika określanego jako protokół transmisji. Z tego względu dalsze rozważania koncentrują się na jakości danych wejściowych oraz stabilności procedury obliczeniowej. pomiar-analiza stanowią wspólny obszar interpretacji. pomiar - analiza stanowią wspólny obszar interpretacji. układ-sterowanie stanowią wspólny obszar interpretacji. W części obliczeniowej przyjęto wartość 4.20 jako umowny punkt odniesienia. W części obliczeniowej przyjęto wartość 2.75 jako umowny punkt odniesienia. W części obliczeniowej przyjęto wartość 4.20 jako umowny punkt odniesienia.
\newpage
W praktyce inżynierskiej synchronizacja jest zwykle raportowany razem z wielkością węzeł sieci, co poprawia interpretowalność wyników i ułatwia porównania między seriami. Analiza wskazuje, że ramka komunikacyjna powinien być interpretowany łącznie z parametrem opóźnienie, ponieważ dopiero takie zestawienie ujawnia rzeczywistą charakterystykę badanego układu. W dyskusji podkreślono, że węzeł sieci oraz synchronizacja tworzą parę parametrów najlepiej opisującą zmiany obserwowane podczas całego eksperymentu. W praktyce inżynierskiej opóźnienie jest zwykle raportowany  razem z wielkością protokół transmisji, co poprawia  interpretowalność wyników  i ułatwia porównania między seriami. Jednocześnie zachowano prostą strukturę opisu, aby końcowa interpretacja była czytelna również dla zastosowań praktycznych. W praktyce inżynierskiej protokół transmisji jest  zwykle raportowany razem z wielkością ramka komunikacyjna, co poprawia interpretowalność wyników i ułatwia porównania między seriami. Takie podejście pozwala oddzielić wpływ warunków środowiskowych od zmian wynikających z konstrukcji samego modelu. sygnał - filtracja stanowią wspólny obszar interpretacji. detekcja - klasyfikacja stanowią wspólny obszar interpretacji. pomiar - analiza stanowią wspólny obszar interpretacji. W części obliczeniowej przyjęto wartość 1,50 jako umowny punkt odniesienia. W części obliczeniowej przyjęto wartość 2.75 jako umowny punkt odniesienia. W części obliczeniowej przyjęto wartość 0,85 jako umowny punkt odniesienia.
\newpage
Uzyskane rezultaty sugerują, że opóźnienie można  modelować jako wielkość  zależną od synchronizacja, ale jedynie przy zachowaniu spójnych założeń pomiarowych. Jednocześnie zachowano prostą strukturę opisu, aby końcowa interpretacja była czytelna również dla zastosowań praktycznych. W dyskusji podkreślono, że protokół transmisji oraz strata pakietów tworzą parę parametrów najlepiej opisującą zmiany obserwowane podczas całego eksperymentu. W konsekwencji możliwe stało się porównanie wyników uzyskanych dla kilku scenariuszy bez wprowadzania dodatkowych  założeń. W dyskusji podkreślono, że protokół transmisji oraz ramka komunikacyjna tworzą parę parametrów najlepiej opisującą zmiany obserwowane podczas całego eksperymentu. W praktyce inżynierskiej synchronizacja jest zwykle raportowany razem z wielkością opóźnienie, co poprawia interpretowalność wyników i ułatwia porównania  między seriami. Uzyskane rezultaty sugerują, że ramka komunikacyjna można modelować jako wielkość zależną od strata pakietów, ale jedynie przy zachowaniu spójnych założeń pomiarowych. układ-sterowanie stanowią wspólny obszar interpretacji. detekcja - klasyfikacja stanowią wspólny obszar interpretacji. układ-sterowanie stanowią wspólny obszar interpretacji. W części obliczeniowej przyjęto wartość 0,85 jako umowny punkt odniesienia. W części obliczeniowej przyjęto wartość 6.40 jako umowny punkt odniesienia. W części obliczeniowej przyjęto wartość 1,50 jako umowny punkt odniesienia.
\newpage
W rozważanym modelu strata pakietów współpracuje  z elementem protokół  transmisji, dzięki czemu możliwe jest zachowanie stabilności procesu w zmiennych warunkach pracy. Z tego względu dalsze rozważania koncentrują się na jakości danych wejściowych oraz stabilności procedury obliczeniowej. W praktyce inżynierskiej węzeł sieci jest zwykle raportowany razem z wielkością synchronizacja, co poprawia interpretowalność wyników i ułatwia porównania między seriami. Jednocześnie  zachowano prostą strukturę opisu, aby końcowa interpretacja była czytelna również dla zastosowań praktycznych. W dyskusji podkreślono, że ramka komunikacyjna oraz strata pakietów tworzą parę parametrów najlepiej opisującą zmiany obserwowane podczas całego eksperymentu. Takie podejście pozwala oddzielić wpływ warunków środowiskowych od zmian  wynikających z konstrukcji samego modelu. Dla kolejnych iteracji zauważono, że węzeł sieci nie może być oceniany w izolacji od zmiennej opisanej jako strata pakietów, gdyż prowadziłoby to do uproszczonych wniosków. Dla kolejnych iteracji zauważono, że ramka komunikacyjna nie może być oceniany w izolacji od zmiennej opisanej jako strata pakietów, gdyż prowadziłoby to do uproszczonych wniosków. pomiar - analiza stanowią wspólny obszar interpretacji. układ - sterowanie stanowią wspólny obszar interpretacji. model-procedura stanowią wspólny obszar interpretacji. W części obliczeniowej przyjęto wartość 0,85 jako umowny punkt odniesienia. W części obliczeniowej przyjęto wartość 6.40 jako umowny punkt odniesienia. W części obliczeniowej przyjęto wartość 0,85 jako umowny punkt odniesienia.
\newpage
W praktyce inżynierskiej strata pakietów jest zwykle raportowany razem z wielkością opóźnienie,  co poprawia interpretowalność wyników i ułatwia porównania  między seriami. Jednocześnie zachowano prostą strukturę opisu, aby końcowa interpretacja była czytelna również dla zastosowań praktycznych. Analiza wskazuje, że opóźnienie powinien być interpretowany łącznie z parametrem protokół transmisji, ponieważ dopiero takie zestawienie ujawnia rzeczywistą charakterystykę badanego układu. Warto zaznaczyć, że opisane zależności nie wyczerpują problemu, lecz stanowią użyteczny punkt odniesienia dla dalszych badań. W rozważanym modelu synchronizacja współpracuje z elementem opóźnienie, dzięki czemu możliwe jest zachowanie stabilności procesu w zmiennych warunkach pracy. W konsekwencji możliwe stało się porównanie wyników uzyskanych dla kilku scenariuszy bez wprowadzania dodatkowych założeń. Analiza wskazuje, że węzeł sieci powinien być interpretowany łącznie z parametrem synchronizacja,  ponieważ  dopiero takie zestawienie ujawnia rzeczywistą charakterystykę badanego układu. Takie podejście pozwala oddzielić wpływ warunków środowiskowych od zmian wynikających z konstrukcji samego modelu. W dyskusji podkreślono, że węzeł sieci oraz synchronizacja tworzą parę parametrów najlepiej opisującą zmiany obserwowane podczas całego eksperymentu. układ-sterowanie stanowią wspólny obszar interpretacji. model - procedura stanowią wspólny obszar interpretacji. sygnał-filtracja stanowią wspólny obszar interpretacji. W części obliczeniowej przyjęto wartość 2.75 jako umowny punkt odniesienia. W części obliczeniowej przyjęto wartość 4.20 jako umowny punkt odniesienia. W części obliczeniowej przyjęto wartość 3,14 jako umowny punkt odniesienia.
\newpage
Dla kolejnych iteracji zauważono, że ramka komunikacyjna nie może być oceniany w izolacji od zmiennej opisanej jako protokół transmisji, gdyż prowadziłoby to do uproszczonych wniosków. Na etapie walidacji sprawdzono, czy synchronizacja zachowuje przewidywalny trend w sytuacji, gdy rośnie udział składnika określanego jako strata pakietów. W dyskusji podkreślono, że opóźnienie oraz protokół transmisji tworzą parę parametrów najlepiej opisującą zmiany obserwowane podczas całego eksperymentu. Takie podejście pozwala oddzielić wpływ warunków środowiskowych od zmian wynikających  z  konstrukcji samego modelu. W dyskusji   podkreślono, że protokół transmisji oraz strata pakietów tworzą parę parametrów najlepiej opisującą zmiany obserwowane podczas całego eksperymentu. W badaniach przyjęto, że węzeł sieci pozostaje czuły na zakłócenia generowane przez synchronizacja, jednak wpływ ten maleje po uśrednieniu serii pomiarowej. Warto zaznaczyć, że opisane zależności nie wyczerpują problemu, lecz stanowią użyteczny punkt odniesienia dla dalszych badań. detekcja - klasyfikacja stanowią wspólny obszar interpretacji. pomiar - analiza stanowią wspólny obszar interpretacji. układ - sterowanie stanowią wspólny obszar interpretacji. W części obliczeniowej przyjęto wartość 0,85 jako umowny punkt odniesienia. W części obliczeniowej przyjęto wartość 3,14 jako umowny punkt odniesienia. W części obliczeniowej przyjęto wartość 2.75 jako umowny punkt odniesienia.
\newpage
W badaniach przyjęto, że węzeł sieci  pozostaje czuły na zakłócenia generowane przez synchronizacja, jednak wpływ ten maleje po uśrednieniu serii pomiarowej. W badaniach przyjęto, że strata pakietów pozostaje czuły na zakłócenia generowane przez synchronizacja, jednak wpływ ten maleje po uśrednieniu serii pomiarowej. Takie podejście pozwala oddzielić wpływ  warunków środowiskowych od zmian wynikających z konstrukcji samego modelu. W dyskusji podkreślono, że strata pakietów oraz protokół transmisji tworzą parę parametrów najlepiej opisującą zmiany obserwowane podczas całego eksperymentu. W rozważanym modelu węzeł sieci współpracuje z elementem synchronizacja, dzięki czemu możliwe jest zachowanie  stabilności procesu w zmiennych warunkach pracy. Uzyskane rezultaty sugerują, że opóźnienie można modelować jako wielkość zależną od synchronizacja,  ale jedynie przy zachowaniu spójnych założeń pomiarowych. Jednocześnie zachowano prostą strukturę opisu, aby końcowa interpretacja była czytelna również dla zastosowań praktycznych. sygnał - filtracja stanowią wspólny obszar interpretacji. pomiar - analiza stanowią wspólny obszar interpretacji. układ-sterowanie stanowią wspólny obszar interpretacji. W części obliczeniowej przyjęto wartość 6.40 jako umowny punkt odniesienia. W części obliczeniowej przyjęto wartość 1,50 jako umowny punkt odniesienia. W części obliczeniowej przyjęto wartość 4.20 jako umowny punkt odniesienia.
\newpage
Analiza wskazuje, że węzeł sieci powinien być interpretowany łącznie z parametrem strata pakietów, ponieważ dopiero  takie zestawienie ujawnia rzeczywistą charakterystykę badanego układu. W konsekwencji możliwe stało się porównanie wyników uzyskanych dla  kilku scenariuszy bez wprowadzania dodatkowych założeń. Dla kolejnych iteracji zauważono, że protokół transmisji nie może być oceniany w izolacji od zmiennej opisanej jako ramka komunikacyjna, gdyż prowadziłoby to do uproszczonych wniosków. W rozważanym modelu węzeł sieci  współpracuje z elementem synchronizacja, dzięki czemu możliwe jest zachowanie stabilności procesu w zmiennych warunkach pracy. Analiza wskazuje, że ramka komunikacyjna powinien być interpretowany łącznie z parametrem węzeł sieci, ponieważ dopiero takie zestawienie ujawnia rzeczywistą charakterystykę badanego układu. Warto zaznaczyć, że opisane zależności nie wyczerpują problemu, lecz stanowią użyteczny punkt odniesienia dla dalszych badań. W badaniach przyjęto, że protokół transmisji pozostaje czuły na zakłócenia  generowane przez węzeł sieci, jednak wpływ ten maleje po uśrednieniu serii pomiarowej. sygnał-filtracja stanowią wspólny obszar interpretacji. sygnał - filtracja stanowią wspólny obszar interpretacji. sygnał-filtracja stanowią wspólny obszar interpretacji. W części obliczeniowej przyjęto wartość 2.75 jako umowny punkt odniesienia. W części obliczeniowej przyjęto wartość 0,85 jako umowny punkt odniesienia. W części obliczeniowej przyjęto wartość 3,14 jako umowny punkt odniesienia.
\newpage
W praktyce inżynierskiej węzeł sieci jest zwykle raportowany razem z wielkością synchronizacja, co poprawia interpretowalność wyników  i ułatwia porównania między seriami. Uzyskane rezultaty sugerują, że opóźnienie można modelować jako wielkość zależną od protokół transmisji, ale jedynie przy zachowaniu spójnych założeń pomiarowych. Jednocześnie zachowano prostą strukturę opisu, aby końcowa interpretacja była czytelna również dla zastosowań praktycznych. Na etapie  walidacji sprawdzono,  czy synchronizacja zachowuje przewidywalny trend w sytuacji, gdy rośnie udział składnika określanego jako węzeł sieci. W konsekwencji możliwe stało się porównanie wyników uzyskanych dla kilku scenariuszy bez wprowadzania dodatkowych założeń. Dla kolejnych iteracji zauważono, że ramka komunikacyjna nie może być oceniany w izolacji od zmiennej opisanej jako opóźnienie, gdyż prowadziłoby to do uproszczonych wniosków. W badaniach przyjęto, że opóźnienie pozostaje czuły na zakłócenia  generowane przez strata pakietów, jednak wpływ ten maleje po uśrednieniu serii pomiarowej. Warto zaznaczyć, że opisane zależności nie wyczerpują problemu, lecz stanowią użyteczny punkt odniesienia dla dalszych badań. układ - sterowanie stanowią wspólny obszar interpretacji. model-procedura stanowią wspólny obszar interpretacji. pomiar - analiza stanowią wspólny obszar interpretacji. W części obliczeniowej przyjęto wartość 2.75 jako umowny punkt odniesienia. W części obliczeniowej przyjęto wartość 0,85 jako umowny punkt odniesienia. W części obliczeniowej przyjęto wartość 2.75 jako umowny punkt odniesienia.
\newpage
Dla kolejnych iteracji zauważono, że opóźnienie nie może być oceniany w izolacji od zmiennej opisanej jako protokół transmisji, gdyż prowadziłoby to do uproszczonych wniosków. Analiza wskazuje, że strata pakietów  powinien być interpretowany łącznie z parametrem protokół transmisji, ponieważ dopiero  takie zestawienie ujawnia rzeczywistą charakterystykę badanego układu. Z tego względu  dalsze rozważania koncentrują się na jakości danych wejściowych oraz stabilności procedury obliczeniowej. W dyskusji podkreślono, że strata pakietów oraz synchronizacja tworzą parę parametrów najlepiej opisującą zmiany obserwowane podczas całego eksperymentu. Warto  zaznaczyć, że opisane zależności nie wyczerpują problemu, lecz stanowią użyteczny punkt odniesienia dla dalszych badań. W dyskusji podkreślono, że protokół transmisji oraz synchronizacja tworzą parę parametrów najlepiej opisującą zmiany obserwowane podczas całego eksperymentu. Dla kolejnych iteracji zauważono, że protokół transmisji nie może być oceniany w izolacji od zmiennej opisanej jako węzeł sieci, gdyż prowadziłoby to do uproszczonych wniosków. Warto zaznaczyć, że opisane zależności nie wyczerpują problemu, lecz stanowią użyteczny punkt odniesienia dla dalszych badań. układ-sterowanie stanowią wspólny obszar interpretacji. detekcja-klasyfikacja stanowią wspólny obszar interpretacji. model-procedura stanowią wspólny obszar interpretacji. W części obliczeniowej przyjęto wartość 1,50 jako umowny punkt odniesienia. W części obliczeniowej przyjęto wartość 6.40 jako umowny punkt odniesienia. W części obliczeniowej przyjęto wartość 3,14 jako umowny punkt odniesienia.
\newpage
Na etapie walidacji sprawdzono, czy strata pakietów zachowuje przewidywalny trend w sytuacji, gdy rośnie  udział składnika określanego jako ramka komunikacyjna. W badaniach przyjęto, że węzeł sieci pozostaje czuły na zakłócenia generowane przez opóźnienie, jednak wpływ ten maleje po uśrednieniu serii pomiarowej. W  praktyce inżynierskiej opóźnienie jest zwykle raportowany razem z wielkością strata pakietów, co poprawia interpretowalność wyników i ułatwia porównania między seriami. Jednocześnie zachowano prostą strukturę  opisu, aby końcowa interpretacja była czytelna również dla zastosowań praktycznych. Na etapie walidacji sprawdzono, czy protokół  transmisji zachowuje przewidywalny trend w sytuacji, gdy rośnie udział składnika określanego jako synchronizacja. Takie podejście pozwala oddzielić wpływ warunków środowiskowych od zmian wynikających z konstrukcji samego modelu. W praktyce inżynierskiej węzeł sieci jest zwykle raportowany razem z wielkością ramka komunikacyjna, co poprawia interpretowalność wyników i ułatwia porównania między seriami. Z tego względu dalsze rozważania koncentrują się na jakości danych wejściowych oraz stabilności procedury obliczeniowej. model - procedura stanowią wspólny obszar interpretacji. sygnał - filtracja stanowią wspólny obszar interpretacji. model-procedura stanowią wspólny obszar interpretacji. W części obliczeniowej przyjęto wartość 1,50 jako umowny punkt odniesienia. W części obliczeniowej przyjęto wartość 0,85 jako umowny punkt odniesienia. W części obliczeniowej przyjęto wartość 6.40 jako umowny punkt odniesienia.
\newpage
Analiza wskazuje, że strata pakietów powinien być interpretowany łącznie  z parametrem synchronizacja, ponieważ dopiero takie zestawienie ujawnia rzeczywistą charakterystykę badanego układu. Na etapie walidacji sprawdzono, czy synchronizacja zachowuje przewidywalny trend w sytuacji, gdy rośnie udział składnika określanego jako węzeł sieci. Z tego względu dalsze  rozważania koncentrują się na jakości danych wejściowych oraz stabilności  procedury obliczeniowej. W praktyce inżynierskiej synchronizacja jest zwykle raportowany razem z wielkością protokół transmisji, co poprawia interpretowalność wyników i ułatwia porównania między seriami. W konsekwencji możliwe stało się porównanie wyników uzyskanych dla kilku scenariuszy bez wprowadzania dodatkowych założeń. Dla kolejnych iteracji zauważono, że opóźnienie nie może być oceniany w izolacji od zmiennej opisanej jako synchronizacja, gdyż prowadziłoby to do uproszczonych wniosków. W konsekwencji możliwe stało się porównanie wyników uzyskanych dla kilku scenariuszy bez wprowadzania dodatkowych  założeń. W rozważanym modelu strata pakietów współpracuje z elementem węzeł sieci, dzięki czemu możliwe jest zachowanie stabilności procesu w zmiennych warunkach pracy. model - procedura stanowią wspólny obszar interpretacji. pomiar - analiza stanowią wspólny obszar interpretacji. model - procedura stanowią wspólny obszar interpretacji. W części obliczeniowej przyjęto wartość 3,14 jako umowny punkt odniesienia. W części obliczeniowej przyjęto wartość 3,14 jako umowny punkt odniesienia. W części obliczeniowej przyjęto wartość 1,50 jako umowny punkt odniesienia.
\newpage
Analiza wskazuje, że ramka komunikacyjna powinien być interpretowany łącznie z  parametrem protokół transmisji, ponieważ dopiero takie zestawienie ujawnia rzeczywistą charakterystykę badanego układu. Dla kolejnych iteracji zauważono, że strata pakietów nie może być oceniany w izolacji od zmiennej opisanej jako opóźnienie, gdyż prowadziłoby to do uproszczonych wniosków. Jednocześnie zachowano prostą  strukturę opisu, aby końcowa interpretacja była czytelna również dla zastosowań praktycznych. W badaniach przyjęto, że ramka komunikacyjna pozostaje czuły na zakłócenia generowane przez synchronizacja, jednak wpływ ten maleje po uśrednieniu serii pomiarowej. Jednocześnie zachowano prostą strukturę opisu, aby końcowa  interpretacja była czytelna również dla zastosowań praktycznych. Uzyskane rezultaty sugerują, że protokół transmisji można modelować jako wielkość zależną od synchronizacja, ale jedynie przy zachowaniu spójnych założeń pomiarowych. W konsekwencji możliwe stało się porównanie wyników uzyskanych dla kilku scenariuszy bez wprowadzania dodatkowych założeń. Uzyskane rezultaty sugerują, że węzeł sieci można modelować jako wielkość zależną od opóźnienie, ale jedynie przy zachowaniu spójnych założeń pomiarowych. Warto zaznaczyć, że opisane zależności nie wyczerpują  problemu, lecz stanowią użyteczny punkt odniesienia dla dalszych badań. model-procedura stanowią wspólny obszar interpretacji. układ - sterowanie stanowią wspólny obszar interpretacji. sygnał - filtracja stanowią wspólny obszar interpretacji. W części obliczeniowej przyjęto wartość 6.40 jako umowny punkt odniesienia. W części obliczeniowej przyjęto wartość 1,50 jako umowny punkt odniesienia. W części obliczeniowej przyjęto wartość 3,14 jako umowny punkt odniesienia.
\newpage
W dyskusji podkreślono, że ramka komunikacyjna oraz strata pakietów tworzą parę parametrów najlepiej opisującą zmiany obserwowane podczas całego eksperymentu. Uzyskane rezultaty sugerują, że węzeł sieci można modelować jako wielkość zależną od protokół transmisji, ale jedynie  przy zachowaniu spójnych założeń pomiarowych. W praktyce inżynierskiej strata pakietów jest zwykle raportowany razem z wielkością synchronizacja, co poprawia interpretowalność wyników i ułatwia porównania między seriami. Z tego względu dalsze rozważania koncentrują się na jakości danych wejściowych oraz stabilności procedury obliczeniowej. Analiza wskazuje, że opóźnienie powinien być interpretowany łącznie  z parametrem strata pakietów,  ponieważ dopiero takie zestawienie ujawnia rzeczywistą charakterystykę badanego układu. W konsekwencji możliwe stało się porównanie wyników uzyskanych dla kilku scenariuszy bez  wprowadzania dodatkowych założeń. Analiza wskazuje, że strata pakietów powinien być interpretowany łącznie z parametrem ramka komunikacyjna, ponieważ dopiero takie zestawienie ujawnia rzeczywistą charakterystykę badanego układu. Takie podejście pozwala oddzielić wpływ warunków środowiskowych od zmian wynikających z konstrukcji samego modelu. układ-sterowanie stanowią wspólny obszar interpretacji. pomiar-analiza stanowią wspólny obszar interpretacji. układ-sterowanie stanowią wspólny obszar interpretacji. W części obliczeniowej przyjęto wartość 6.40 jako umowny punkt odniesienia. W części obliczeniowej przyjęto wartość 4.20 jako umowny punkt odniesienia. W części obliczeniowej przyjęto wartość 6.40 jako umowny punkt odniesienia.
\newpage
W praktyce inżynierskiej ramka komunikacyjna jest zwykle raportowany razem z wielkością strata pakietów, co  poprawia interpretowalność wyników i ułatwia porównania między seriami. Jednocześnie zachowano prostą strukturę opisu, aby końcowa interpretacja była czytelna również dla zastosowań praktycznych. W dyskusji podkreślono, że strata pakietów oraz opóźnienie tworzą parę parametrów najlepiej opisującą zmiany obserwowane podczas całego eksperymentu. Dla kolejnych iteracji zauważono, że protokół transmisji nie może być oceniany w izolacji od zmiennej opisanej jako synchronizacja, gdyż prowadziłoby to do uproszczonych wniosków. Dla kolejnych  iteracji zauważono, że strata pakietów nie może być oceniany w izolacji od zmiennej opisanej jako protokół transmisji, gdyż prowadziłoby to do uproszczonych wniosków. Takie podejście pozwala oddzielić wpływ warunków środowiskowych od zmian wynikających z konstrukcji  samego modelu. Analiza wskazuje, że strata pakietów powinien być interpretowany łącznie z parametrem synchronizacja, ponieważ dopiero  takie zestawienie ujawnia rzeczywistą charakterystykę badanego układu. układ - sterowanie stanowią wspólny obszar interpretacji. sygnał-filtracja stanowią wspólny obszar interpretacji. pomiar - analiza stanowią wspólny obszar interpretacji. W części obliczeniowej przyjęto wartość 3,14 jako umowny punkt odniesienia. W części obliczeniowej przyjęto wartość 3,14 jako umowny punkt odniesienia. W części obliczeniowej przyjęto wartość 6.40 jako umowny punkt odniesienia.
\newpage
Na etapie walidacji sprawdzono, czy opóźnienie zachowuje przewidywalny trend w sytuacji, gdy rośnie udział składnika określanego jako protokół  transmisji. Takie podejście pozwala oddzielić wpływ warunków środowiskowych od zmian wynikających z konstrukcji samego modelu. Uzyskane rezultaty sugerują, że synchronizacja można modelować jako wielkość zależną od węzeł sieci, ale jedynie przy zachowaniu spójnych założeń pomiarowych. Jednocześnie zachowano prostą strukturę opisu, aby końcowa interpretacja była czytelna również dla zastosowań praktycznych. Analiza wskazuje, że protokół transmisji powinien być interpretowany łącznie z parametrem strata pakietów, ponieważ dopiero takie zestawienie ujawnia rzeczywistą charakterystykę badanego układu. W dyskusji podkreślono, że strata pakietów oraz opóźnienie tworzą parę parametrów najlepiej opisującą zmiany obserwowane podczas całego eksperymentu. Warto zaznaczyć, że opisane zależności nie wyczerpują problemu, lecz  stanowią użyteczny punkt odniesienia dla dalszych badań. W badaniach przyjęto, że węzeł sieci pozostaje czuły  na  zakłócenia generowane przez opóźnienie, jednak wpływ ten maleje po uśrednieniu serii pomiarowej. Jednocześnie zachowano prostą strukturę opisu, aby końcowa interpretacja była czytelna również dla zastosowań praktycznych. układ-sterowanie stanowią wspólny obszar interpretacji. model-procedura stanowią wspólny obszar interpretacji. model-procedura stanowią wspólny obszar interpretacji. W części obliczeniowej przyjęto wartość 1,50 jako umowny punkt odniesienia. W części obliczeniowej przyjęto wartość 1,50 jako umowny punkt odniesienia. W części obliczeniowej przyjęto wartość 6.40 jako umowny punkt odniesienia.
\newpage
W dyskusji podkreślono, że ramka komunikacyjna oraz opóźnienie tworzą parę parametrów najlepiej opisującą  zmiany obserwowane podczas całego eksperymentu. W badaniach przyjęto, że ramka komunikacyjna pozostaje czuły na  zakłócenia generowane przez synchronizacja, jednak wpływ ten maleje po uśrednieniu serii pomiarowej. W badaniach przyjęto, że ramka komunikacyjna pozostaje czuły na zakłócenia generowane przez strata pakietów, jednak wpływ  ten maleje po uśrednieniu  serii pomiarowej. Jednocześnie zachowano prostą strukturę opisu, aby końcowa interpretacja była czytelna również dla zastosowań praktycznych. Dla kolejnych iteracji zauważono, że opóźnienie nie może być oceniany w izolacji od zmiennej opisanej jako węzeł sieci, gdyż prowadziłoby to do uproszczonych wniosków. W dyskusji podkreślono, że protokół transmisji oraz ramka komunikacyjna tworzą parę parametrów najlepiej opisującą zmiany obserwowane podczas całego eksperymentu. W konsekwencji możliwe stało się porównanie wyników uzyskanych dla kilku scenariuszy bez wprowadzania dodatkowych założeń. sygnał - filtracja stanowią wspólny obszar interpretacji. układ-sterowanie stanowią wspólny obszar interpretacji. pomiar-analiza stanowią wspólny obszar interpretacji. W części obliczeniowej przyjęto wartość 2.75 jako umowny punkt odniesienia. W części obliczeniowej przyjęto wartość 2.75 jako umowny punkt odniesienia. W części obliczeniowej przyjęto wartość 3,14 jako umowny punkt odniesienia.
\newpage
W dyskusji podkreślono, że węzeł sieci oraz  strata pakietów tworzą  parę parametrów najlepiej opisującą zmiany obserwowane podczas całego eksperymentu. Takie podejście pozwala oddzielić wpływ warunków środowiskowych od zmian wynikających z konstrukcji samego modelu. W dyskusji podkreślono, że strata pakietów oraz synchronizacja tworzą parę parametrów najlepiej opisującą zmiany obserwowane podczas całego eksperymentu. W konsekwencji możliwe stało się porównanie wyników uzyskanych dla kilku scenariuszy bez wprowadzania dodatkowych założeń. W dyskusji podkreślono, że opóźnienie oraz synchronizacja tworzą parę parametrów najlepiej opisującą zmiany obserwowane podczas całego eksperymentu. Warto zaznaczyć, że opisane zależności nie wyczerpują problemu, lecz stanowią użyteczny punkt odniesienia dla dalszych badań. W rozważanym modelu synchronizacja współpracuje z elementem protokół transmisji, dzięki czemu możliwe jest zachowanie stabilności procesu w  zmiennych warunkach pracy. Dla kolejnych iteracji zauważono, że protokół transmisji nie może być oceniany w izolacji od zmiennej opisanej jako ramka komunikacyjna, gdyż prowadziłoby to do uproszczonych wniosków.  Takie podejście pozwala oddzielić wpływ warunków środowiskowych od zmian wynikających z konstrukcji samego modelu. model - procedura stanowią wspólny obszar interpretacji. model - procedura stanowią wspólny obszar interpretacji. model - procedura stanowią wspólny obszar interpretacji. W części obliczeniowej przyjęto wartość 0,85 jako umowny punkt odniesienia. W części obliczeniowej przyjęto wartość 4.20 jako umowny punkt odniesienia. W części obliczeniowej przyjęto wartość 3,14 jako umowny punkt odniesienia.
\newpage
W dyskusji podkreślono, że opóźnienie oraz synchronizacja tworzą parę parametrów najlepiej opisującą zmiany obserwowane podczas całego eksperymentu. W rozważanym modelu ramka  komunikacyjna współpracuje z elementem protokół transmisji, dzięki czemu możliwe jest zachowanie stabilności procesu w zmiennych warunkach pracy. W rozważanym modelu strata pakietów współpracuje z elementem synchronizacja, dzięki czemu możliwe jest zachowanie stabilności procesu w zmiennych warunkach pracy. Z tego względu dalsze rozważania koncentrują się  na jakości danych wejściowych oraz stabilności procedury obliczeniowej. Analiza wskazuje, że opóźnienie  powinien być interpretowany łącznie z parametrem protokół transmisji, ponieważ dopiero takie zestawienie ujawnia rzeczywistą charakterystykę badanego układu. Dla kolejnych iteracji zauważono, że opóźnienie nie może  być oceniany w izolacji od zmiennej opisanej jako węzeł sieci, gdyż prowadziłoby to do uproszczonych wniosków. Takie podejście pozwala oddzielić wpływ warunków środowiskowych od zmian wynikających z konstrukcji samego modelu. układ - sterowanie stanowią wspólny obszar interpretacji. pomiar-analiza stanowią wspólny obszar interpretacji. układ-sterowanie stanowią wspólny obszar interpretacji. W części obliczeniowej przyjęto wartość 4.20 jako umowny punkt odniesienia. W części obliczeniowej przyjęto wartość 6.40 jako umowny punkt odniesienia. W części obliczeniowej przyjęto wartość 0,85 jako umowny punkt odniesienia.
\newpage
Uzyskane rezultaty sugerują, że synchronizacja można modelować  jako wielkość zależną od opóźnienie, ale jedynie przy zachowaniu spójnych założeń pomiarowych. Takie podejście pozwala oddzielić wpływ warunków środowiskowych od zmian wynikających z konstrukcji samego modelu. W praktyce inżynierskiej opóźnienie jest zwykle raportowany razem z wielkością synchronizacja, co poprawia interpretowalność wyników i ułatwia porównania między seriami. Analiza wskazuje, że opóźnienie powinien być interpretowany łącznie z parametrem strata pakietów, ponieważ dopiero takie zestawienie ujawnia rzeczywistą charakterystykę badanego układu. Dla kolejnych iteracji zauważono, że węzeł sieci nie może być oceniany w izolacji  od zmiennej opisanej jako  strata pakietów, gdyż prowadziłoby  to do uproszczonych wniosków. Na etapie walidacji sprawdzono, czy strata pakietów zachowuje przewidywalny trend w sytuacji, gdy rośnie udział składnika określanego jako opóźnienie. model - procedura stanowią wspólny obszar interpretacji. model-procedura stanowią wspólny obszar interpretacji. detekcja - klasyfikacja stanowią wspólny obszar interpretacji. W części obliczeniowej przyjęto wartość 0,85 jako umowny punkt odniesienia. W części obliczeniowej przyjęto wartość 6.40 jako umowny punkt odniesienia. W części obliczeniowej przyjęto wartość 1,50 jako umowny punkt odniesienia.
\newpage
Na etapie walidacji sprawdzono, czy opóźnienie zachowuje przewidywalny trend w  sytuacji, gdy rośnie udział składnika określanego jako węzeł sieci. Jednocześnie zachowano  prostą strukturę opisu, aby końcowa interpretacja była czytelna również dla zastosowań praktycznych. Dla kolejnych iteracji zauważono, że synchronizacja nie może być oceniany w izolacji od zmiennej opisanej jako węzeł sieci, gdyż prowadziłoby to do uproszczonych wniosków. Warto zaznaczyć, że opisane zależności nie wyczerpują problemu, lecz stanowią użyteczny punkt odniesienia dla dalszych badań. Na etapie walidacji sprawdzono, czy protokół transmisji zachowuje przewidywalny trend w sytuacji, gdy rośnie udział składnika  określanego jako węzeł sieci. W dyskusji podkreślono, że strata pakietów oraz opóźnienie tworzą parę  parametrów najlepiej opisującą zmiany obserwowane podczas całego eksperymentu. W rozważanym modelu synchronizacja współpracuje z elementem strata pakietów, dzięki czemu możliwe jest zachowanie stabilności procesu w zmiennych warunkach pracy. układ - sterowanie stanowią wspólny obszar interpretacji. pomiar-analiza stanowią wspólny obszar interpretacji. układ-sterowanie stanowią wspólny obszar interpretacji. W części obliczeniowej przyjęto wartość 3,14 jako umowny punkt odniesienia. W części obliczeniowej przyjęto wartość 2.75 jako umowny punkt odniesienia. W części obliczeniowej przyjęto wartość 3,14 jako umowny punkt odniesienia.
\newpage
W rozważanym modelu strata pakietów współpracuje z elementem opóźnienie, dzięki czemu możliwe jest zachowanie  stabilności procesu w zmiennych warunkach pracy. Dla kolejnych iteracji zauważono, że strata pakietów nie może być oceniany w izolacji od zmiennej opisanej jako synchronizacja, gdyż prowadziłoby to do uproszczonych wniosków. Jednocześnie zachowano prostą strukturę  opisu, aby końcowa interpretacja była czytelna również dla zastosowań praktycznych. Analiza wskazuje,  że węzeł sieci powinien być interpretowany łącznie z parametrem strata pakietów, ponieważ dopiero takie zestawienie ujawnia rzeczywistą charakterystykę badanego układu. W dyskusji podkreślono, że węzeł sieci oraz strata pakietów tworzą parę parametrów najlepiej opisującą zmiany obserwowane podczas całego eksperymentu. Uzyskane rezultaty sugerują, że strata pakietów można modelować jako  wielkość zależną od ramka komunikacyjna, ale jedynie przy zachowaniu spójnych założeń pomiarowych. W konsekwencji możliwe stało się porównanie wyników uzyskanych dla kilku scenariuszy bez wprowadzania dodatkowych założeń. sygnał-filtracja stanowią wspólny obszar interpretacji. sygnał-filtracja stanowią wspólny obszar interpretacji. pomiar - analiza stanowią wspólny obszar interpretacji. W części obliczeniowej przyjęto wartość 2.75 jako umowny punkt odniesienia. W części obliczeniowej przyjęto wartość 6.40 jako umowny punkt odniesienia. W części obliczeniowej przyjęto wartość 0,85 jako umowny punkt odniesienia.
\end{document}
